\documentclass[12pt]{article}
\usepackage[T1]{fontenc}
\usepackage[utf8]{inputenc}
\usepackage{lmodern}
\usepackage{textcomp}
\usepackage{enumitem}
\usepackage{geometry}
\geometry{margin=1in}
\begin{document}
\begin{center}
Answer Key - Chemistry - Undergraduate Level Assessment 1 \
\small (Answers are ordered exactly as in the question paper.)
\end{center}

\section*{Multiple Choice}
\begin{enumerate}[label=\arabic*.]
\item \textbf{Answer:} D
\\ \textit{Options:} A) Melting point; B) Boiling point; C) Polarizability; D) First ionization energy
\\ \textit{Note:} The first ionization energy of argon is lower than that of neon because argon has a larger atomic radius and more electron shielding, which reduces the effective nuclear charge experienced by its outermost electrons, making them easier to remove.
\item \textbf{Answer:} D
\\ \textit{Options:} A) 10.8 x 10^-5; B) 4.11 x 10^-5; C) 3.43 x 10^-5; D) 1.71 x 10^-5
\\ \textit{Note:} The correct answer of 1.71 x 10^-5 for the equilibrium polarization of 13 C nuclei in a 20.0 T magnetic field at 300 K is derived from the Boltzmann distribution, which predicts the population difference between nuclear spin states that results in a net magnetization proportional to the magnetic field strength and inversely related to temperature.
\item \textbf{Answer:} D
\\ \textit{Options:} A) Q = 1012; B) Q = 2012; C) Q = 3012; D) Q = 6012
\\ \textit{Note:} The Q-factor, defined as the ratio of the resonant frequency to the bandwidth, indicates the energy loss relative to the stored energy in the cavity, and for an X-band EPR cavity with a bandwidth of 1.58 MHz, a Q-factor of 6012 suggests a high level of resonance and low energy loss.
\item \textbf{Answer:} A
\\ \textit{Options:} A) 3; B) 5; C) 8; D) 4
\\ \textit{Note:} The correct answer is 3 because the 55 Mn nuclide has a nuclear spin of 5/2, which corresponds to three allowed energy levels based on the possible projections of the nuclear spin in a magnetic field.
\item \textbf{Answer:} B
\\ \textit{Options:} A) The I- : IO 3- ratio is 3:1.; B) The MnO 4- : I- ratio is 6:5.; C) The MnO 4- : Mn$_{2}$+ ratio is 3:1.; D) The H+ : I- ratio is 2:1.
\\ \textit{Note:} The correct answer is based on the stoichiometry of the balanced redox reaction, where 6 electrons are transferred from 5 iodide ions (I-) to one permanganate ion (MnO 4-), resulting in a ratio of 6:5 for MnO 4- to I-.
\end{enumerate}

\section*{True / False}
\begin{enumerate}[label=\arabic*.]
\setcounter{enumi}{5}
\item \textbf{Answer:} False
\\ \textit{Options:} A) True; B) False
\\ \textit{Note:} The statement is false because, despite potassium (K) being a larger atom than scandium (Sc) in the periodic table, the ionic radius of K+ is larger than that of Sc$_{3}$+ due to the greater effective nuclear charge and loss of three electrons in Sc$_{3}$+, which results in a stronger attraction of the remaining electrons towards the nucleus, leading to a smaller ionic radius.
\item \textbf{Answer:} False
\\ \textit{Options:} A) True; B) False
\\ \textit{Note:} Infrared spectroscopy is not a light-scattering technique; it involves the absorption of infrared light to analyze molecular vibrations.
\item \textbf{Answer:} False
\\ \textit{Options:} A) True; B) False
\\ \textit{Note:} The statement is false because NH 3 (ammonia) is actually a weak base in liquid ammonia, and it does not have a high tendency to donate protons; rather, it tends to accept protons, making it a Lewis base rather than a strong proton donor.
\item \textbf{Answer:} True
\\ \textit{Options:} A) True; B) False
\\ \textit{Note:} Silicon primarily bonds with oxygen in the Earth's crust, forming silicon dioxide (silica) and various silicate minerals, which account for the majority of silicon found in nature.
\item \textbf{Answer:} False
\\ \textit{Options:} A) True; B) False
\\ \textit{Note:} Doubling the volume of a buffer solution by adding water dilutes the concentrations of the buffering agents but does not significantly change the pH because buffers are designed to resist pH changes within a certain range.
\end{enumerate}

\section*{Long Form Response}
\begin{enumerate}[label=\arabic*.]
\setcounter{enumi}{10}
\item \textbf{Answer:} The differing relaxation times, T$_{1}$ and T$_{2}$, are influenced by various factors, particularly the molecular motions occurring at the Larmor frequency. T$_{1}$, or longitudinal relaxation time, is primarily associated with the transfer of energy from the nucleus to its surrounding lattice, which is less sensitive to rapid molecular motions. In contrast, T$_{2}$, or transverse relaxation time, is significantly affected by these molecular motions because it involves the dephasing of nuclear spins due to interactions with other spins and fluctuating magnetic fields at the Larmor frequency. As molecular motions increase, they can disrupt the coherence of the nuclear spins more rapidly, leading to a shorter T$_{2}$ time. Thus, while T$_{1}$ reflects energy exchange processes, T$_{2}$ highlights the impact of molecular dynamics on spin coherence.
\\ \textit{Note:} The answer correctly identifies that T$_{1}$ relaxation is primarily influenced by energy transfer to the surrounding lattice, which is less affected by rapid molecular motions, while T$_{2}$ relaxation is more sensitive to these motions due to the dephasing of nuclear spins at the Larmor frequency.
\item \textbf{Answer:} Ionic strength is a measure of the total concentration of ions in a solution, which influences various chemical properties, including solubility and reaction rates. A$_{0}$.050 M solution of AlCl 3 exhibits higher ionic strength compared to other ionic solutions due to its complete dissociation into three ions: one aluminum ion ($Al^3$+) and three chloride ions (Cl^-). This results in a total effective concentration of ions of 0.050 M AlCl 3 contributing to 0.050 M + 0.150 M = 0.200 M ionic strength. In contrast, a 0.050 M solution of a salt like NaCl only contributes 0.100 M to the ionic strength. Therefore, the higher number of ions produced from AlCl 3 significantly increases the ionic strength of the solution.
\\ \textit{Note:} correct because it accurately calculates the ionic strength of a 0.050 M solution of AlCl 3 by considering its complete dissociation into four total ions, resulting in a higher effective concentration of ions compared to other ionic solutions.
\end{enumerate}

\vfill
\noindent\textit{End of Answer Key}
\end{document}